% EDITORIAL QUESTION: Como as federações internacionais e nacionais dividem competências normativas e quais ferramentas científicas reduzem a subjetividade do route-setting?
\section{Normas Oficiais, Governança e Redução da Subjetividade}

A organização da escalada divide-se entre entidades voltadas ao montanhismo em rocha natural e federações dedicadas ao esporte de competição em estruturas artificiais.

\subsection*{1.1 Hierarquia Internacional: UIAA, World Climbing (ex-IFSC), COI e CPI}

% SOURCE [VERIFIED_PRIMARY]: UIAA History (https://www.theuiaa.org/inside-the-uiaa/history/); World Climbing About & Rebrand 10/12/2025 (https://www.worldclimbing.com/ifsc, https://www.worldclimbing.com/news/world-climbing-launches-new-global-identity).
\begin{itemize}
    \item \textbf{UIAA}: Fundada em agosto de 1932 em Chamonix, atua na certificação de segurança de equipamentos (\textit{UIAA Safety Label}), ética de acesso em montanha, preservação ambiental e competições de escalada no gelo \cite{uiaa_history}. Administrou as Copas do Mundo de escalada esportiva de 1989 a 2006 \cite{worldclimbing_about}.
    \item \textbf{World Climbing (antiga IFSC)}: Fundada em 27 de janeiro de 2007 em Frankfurt como \textit{International Federation of Sport Climbing} (IFSC) para viabilizar a autonomia olímpica da modalidade. Em \textbf{10 de dezembro de 2025}, adotou oficialmente o nome \textbf{World Climbing} e rebatizou o circuito mundial para \textit{World Climbing Series} a partir de 2026 \cite{worldclimbing_about, worldclimbing_rebrand}.
    \item \textbf{COI e CPI}: Reconhecida pelo Comitê Olímpico Internacional (COI) desde 2010 e pelo Comitê Paralímpico Internacional (CPI). Estreou em Tóquio 2020 (formato combinado), dividiu-se em duas disputas em Paris 2024 e terá medalhas independentes para as três disciplinas em Los Angeles 2028, além da estreia paralímpica da \textbf{Paraescalada} \cite{worldclimbing_about}.
\end{itemize}

\subsection*{1.2 As Quatro Disciplinas da World Climbing e o Setting de Competição}

% SOURCE [VERIFIED_PRIMARY]: World Climbing / IFSC Rules 2025 (https://www.worldclimbing.com/commissions/routesetters) & Paraclimbing Classification Rules.
A World Climbing regulamenta quatro categorias oficiais \cite{worldclimbing_about, ifsc_rules}:
\begin{enumerate}
    \item \textbf{Boulder}: Blocos curtos em paredes de até 4,5\,m, sem corda, sobre tatames contínuos, pontuando por \textit{Tops}, \textit{Zones} e número de tentativas em janelas de tempo fixas.
    \item \textbf{Guiada (\textit{Lead})}: Vias de no mínimo 15\,m em tentativa única (\textit{onsight}) com limite de 6 minutos, costurando proteções fixas.
    \item \textbf{Velocidade (\textit{Speed})}: Duelo simultâneo em parede padronizada de 15\,m a $5^\circ$ negativos com agarras e posições idênticas globalmente.
    \item \textbf{Paraescalada (\textit{Para Climbing})}: Disputas em \textit{top-rope} divididas por classes funcionais (B1--B3, AU/AL, RP1--RP3) \cite{ifsc_paraclimbing}.
\end{enumerate}

Diferentemente dos ginásios comerciais, \textbf{competições oficiais da World Climbing não atribuem grau numérico público às vias} \cite{ifsc_rules}. A Comissão de \textit{Setters} designa um \textit{Head Setter} certificado cujo objetivo é produzir \textbf{separação estatística} no pódio, equilibrando estilos distintos (coordenação dinâmica, força em negativo, equilíbrio em placa e complexidade mista) com largada obrigatória em quatro fitas \cite{ifsc_rules}.

\subsection*{1.3 Cenário Institucional no Brasil: CBEscalada vs. CBME}

% SOURCE [VERIFIED_PRIMARY]: CBEscalada Sobre & Filiações 2026 (https://cbescalada.org.br/entidade/9/sobre, https://cbescalada.org.br/atletas/74/filiacoes); CBME DT-2016/01 p. 1.
\begin{itemize}
    \item \textbf{CBEscalada}: Fundada em 2014 (originalmente ABEE), é a associação civil que detém a gestão exclusiva da escalada esportiva de competição no Brasil, integrando o Sistema Nacional do Desporto (Lei nº 14.597/2023) \cite{cbescalada_sobre}. Filiada à World Climbing e reconhecida pelo COB e CPB, opera o Centro de Treinamento Nacional em Curitiba e homologa pódios nacionais junto ao Ministério do Esporte para o programa \textbf{Bolsa Atleta} e editais de fomento como o \textit{Projeto Técnicos de Base} \cite{cbescalada_sobre, cbescalada_filiacoes}.
    \item \textbf{CBME}: Entidade máxima do montanhismo e escalada em rocha no Brasil, filiada à UIAA, que congrega federações estaduais e associações regionais (como FEMERJ, FEMESP, FGM, FEPAM, FEMEMG, FEMECE, ACE, AEP, APE e AERN) \cite{cbme_dt01}. Regula diretrizes técnicas de graduação em rocha (DT-2016/01), ética de conquista (DT-2016/10) e manejo ambiental \cite{cbme_dt01, cbme_dt10}.
\end{itemize}

\subsection*{1.4 A Norma CBME DT-2016/01 e o Grau ``À Vista'' vs. ``Grau Trabalhado''}

% SOURCE [VERIFIED_PRIMARY]: CBME DT-2016/01 (03/12/2016), págs. 1, 2 e 5 (https://cbme.org.br/novo/wp-content/uploads/2025/03/CBME-DT-2016_01-Sistema-Brasileiro-de-graduacao-de-vias-de-escalada.pdf).
A \textbf{Diretriz Técnica CBME DT-2016/01} estabelece o sistema brasileiro de graduação em rocha \cite{cbme_dt01}. O quadro abaixo reproduz três trechos determinantes do documento original:

\facsimilebox{EXCERTO DOCUMENTAL VERIFICADO --- CBME: DT-2016/01 (03/12/2016, Páginas 1, 2 e 5)}{%
\textbf{[Página 1 --- Cabeçalho e Entidades Filiadas]}: ``CONFEDERAÇÃO BRASILEIRA DE MONTANHISMO E ESCALADA --- CBME. Documento: CBME: DT-2016/01. Data criação: 03 de dezembro de 2016. Nota: Substitui o sistema de graduação aprovado por Assembleia Geral em 25/08/2007. Entidades filiadas: FGM, FEPAM, FEMESP, FEMERJ, FEMEMG, FEMECE e as Associações ACE, AEP, APE e Associação de Escaladores do Rio Grande do Norte (AERN).''\\[4pt]
\textbf{[Página 2 --- Seção 2. INTRODUÇÃO]}: ``Uma das vantagens do sistema brasileiro é a menção dos graus geral e do lance mais difícil da via em separado (...). \textbf{Lembramos que na atribuição do grau a uma via considera-se que o escalador está guiando e escalando `à vista', isto é, sem conhecimento prévio da via.}''\\[4pt]
\textbf{[Página 5 --- Seção 3.3. Vias de Uma Enfiada, Falésias e Boulders]}: ``Para estas vias não há sentido em se atribuir um grau geral e um grau para o lance mais difícil, uma vez que são vias curtas, de comprimento máximo de 50 ou 60 metros. \textbf{Então o grau geral é abolido, e utiliza-se somente o grau do lance ou sequência mais difícil, em romanos, para expressar a sua dificuldade.} (...) Esta graduação é válida para boulders, falésias e vias curtas em geral.''
}

\subsubsection*{Grau Geral, Crux e o Efeito do Grau Trabalhado (``Decotar'' a Via)}

% SOURCE [INTERPRETATION]: Análise cultural derivada de CBME DT-2016/01 (págs. 2--5, 13) e convenções de boulder em Hueco/Fontainebleau.
\begin{enumerate}
    \item \textbf{Parede Tradicional vs. Boulder}: Em vias longas (\textit{multi-pitch}), o \textbf{Grau Geral} (arábicos ordinais) expressa o comprometimento contínuo e o \textbf{Grau do Crux} (romanos com \textit{sup} até VIsup e letras \textit{a, b, c} a partir de VIIa) indica o passo mais difícil \cite{cbme_dt01}. Em falésias curtas e boulders, a Seção 3.3 (pág. 5) abole o Grau Geral, adotando apenas o crux (ou a Escala $V$ na prática de campo brasileira) \cite{cbme_dt01, guia_cocalzinho, guia_igatu}.
    \item \textbf{À Vista (\textit{Onsight}) vs. Grau Trabalhado (``Decotado'')}: A CBME fixa o grau sob a premissa da escalada \hlkey{``à vista'' (\textit{onsight})} \cite{cbme_dt01}. Quando o escalador ensaia a sequência repetidamente, memoriza micro-apoios e automatiza o gesto motor até encadeá-la sem hesitação, pratica o \hlwork{\textbf{grau trabalhado} (ou ``decotar'' a via, do francês \textit{décoter})}.
    \item \textbf{Choque Cultural em Ginásios}: Enquanto a parede tradicional cobra o erro de leitura sob exposição, o boulder moderno internacional gradua assumindo o \hlwork{\textbf{grau trabalhado / projeto (\textit{redpoint})}}. Em ginásios, um bloco baseado apenas em um truque de leitura (\textit{beta} oculto) parece dois graus mais duro na primeira tentativa e despenca de dificuldade após ser decotado — exigindo que os \textit{setters} não superestimem graus baseados em enigmas visuais.
\end{enumerate}

\subsection*{1.5 Ferramentas Científicas para Setters: Escala IRCRA}

% SOURCE [VERIFIED_PRIMARY]: Draper et al. (2016), Sports Technology 8(3-4), págs. 85--93 (https://ircra.rocks/reporting-grades-in-climbing-research/).
Existem instrumentos científicos criados para medir e comparar graduações de escalada sem depender da percepção isolada de um único avaliador. A \textbf{IRCRA} desenvolveu uma escala numérica inteira universal (1 a 32, posteriormente expandida até 39 na literatura) para viabilizar análises estatísticas entre sistemas incompatíveis como YDS, Font e $V$-Scale \cite{ircra_draper, pmc_ric_study}.

Para \textit{setters} e gestores de ginásios, embora a IRCRA nunca vá impressa nas etiquetas da parede, ela oferece um diagnóstico prático decisivo: \textbf{revela a não linearidade oculta da Escala $V$} \cite{ircra_draper, youthclimbing_conv}. Entre VB (IRCRA 9) e V1 (IRCRA 14) existem cinco degraus inteiros de exigência motora, enquanto de V1 a V3 avança-se apenas um degrau IRCRA por nível (detalhamento formal e limites científicos no Apêndice~B).

\subsection*{1.6 Redução da Subjetividade: Matriz RIC, Ergonomia e Forerunning}

% SOURCE [VERIFIED_PRIMARY]: PMC11881084 (2024), págs. 2--5 (https://pmc.ncbi.nlm.nih.gov/articles/PMC11881084/); Anderson (2014), caps. 3 e 5; CWA (2023).
Para reduzir a dependência do biotipo, protocolos técnicos de \textit{route-setting} utilizam três barreiras de controle \cite{cwa_setting, anderson_routesetting, pmc_ric_study}:
\begin{enumerate}
    \item \textbf{Matriz RIC (Risco, Intensidade e Complexidade)}: Proposta por Godoffe e Katiyo e formalizada em estudos recentes \cite{godoffe_routesetting, pmc_ric_study}, atribui notas de 1 a 5 a três eixos independentes: \textbf{Risco ($R$)} (comprometimento espacial/queda), \textbf{Intensidade ($I$)} (carga física/força de dedos) e \textbf{Complexidade ($C$)} (leitura e coordenação). Isso garante que estilos opostos (placa técnica vs. teto de força) convirjam para o mesmo grau nominal \cite{agility_ric}.
    \item \textbf{Ergonomia e Controle Morfológico}: Conforme Anderson \cite{anderson_routesetting}, evita-se fabricar dificuldade por meio de posições lesivas (como arqueamento fechado em graus iniciais) e neutraliza-se o viés de envergadura com apoios intermediários de pé ou agarras duplas \cite{paraclimbing_guide}.
    \item \textbf{Forerunning Coletivo e Consenso}: A CWA estabelece que nenhum bloco seja graduado por um único \textit{setter} \cite{cwa_setting}. A testagem prévia por múltiplos escaladores de estaturas distintas corrige quebras de sequência (\textit{beta breaks}), complementada por auditoria estatística via aplicativos digitais após repetições do público \cite{rockrealms_home, topclimbing_platform}.
\end{enumerate}

\noindent\textit{Síntese:} A normatização institucional (World Climbing, CBME, CBEscalada) delimita o escopo de cada escala, enquanto a combinação de IRCRA (diagnóstico linear), Matriz RIC (equilíbrio de estilos) e \textit{forerunning} coletivo transforma a graduação de opinião individual em processo auditável.
